\section{粗粒度 KL 的信息解释}
\label{inf:identification}

\begin{theorem}[桶质量信息下的 KL 下包络]
\label{inf:lower-envelope}
设 $P(v)>0$，每个已知桶质量 $Q_C>0$。仅在定义~\ref{inf:compatible} 的桶质量信息类 $\mathfrak Q_{\mathcal C}$ 上，
\begin{equation}
\inf_{\widetilde Q\in\mathfrak Q_{\mathcal C}}\KL(P\Vert\widetilde Q)
=\sum_CP_C\log\frac{P_C}{Q_C}=K_{\mathcal C}.
\end{equation}
唯一达到下包络的参考分布为 $\widetilde Q(v)=Q_CP^C(v)$，$v\in C$。$\KL(Q\Vert P)$ 同样满足
\begin{equation}
\inf_{\widetilde Q\in\mathfrak Q_{\mathcal C}}\KL(\widetilde Q\Vert P)
=\sum_CQ_C\log\frac{Q_C}{P_C},
\end{equation}
并由同一个分布达到。
\end{theorem}
\begin{proof}
由双向链式分解，
\begin{align}
\KL(P\Vert\widetilde Q)&=K_{\mathcal C}+\sum_CP_C\KL(P^C\Vert\widetilde Q^C),\\
\KL(\widetilde Q\Vert P)&=K^{\mathrm F}_{\mathcal C}+\sum_CQ_C\KL(\widetilde Q^C\Vert P^C).
\end{align}
全部条件 KL 非负，因而两式分别不小于相应粗粒度目标。取 $\widetilde Q^C=P^C$ 时，构造出的 $\widetilde Q(v)=Q_CP^C(v)>0$ 满足全部桶质量约束，两式的条件 KL 同时为零。反之，正权重下条件 KL 和为零要求每个条件 KL 为零，故每个桶内分布必须相同，证明唯一性。
\end{proof}

头部 singleton、尾部单桶时，上式就是 Tail-aware KL。它将未知尾部形状差异设为零，得到桶质量信息能够保证的最小完整 KL。实际 Top-$J$ 信息还具有排名约束，因此这个锐下包络结论不能直接搬到全部真实稀疏 payload。

例如 $P=(0.1,0.8,0.1)$，已知 $Q_1=0.4$、$Q_2+Q_3=0.6$，且 $S=\{1\}$ 是真实 Top-1。仅使用质量约束的最小点为
\begin{equation}
\widetilde Q_2=0.6\frac{0.8}{0.9}=\frac8{15}>0.4.
\end{equation}
它违反 Top-1 排名。加入排名后，Tail-aware 仍是下界，但最小值严格提高：可兼容 $Q_2\leq0.4$，在闭边界最优点 $Q_2=0.4,Q_3=0.2$ 下，尾部条件分布为 $(2/3,1/3)$，与 $P^\tail=(8/9,1/9)$ 不同，条件 KL 严格为正。额外 sampled-action 概率也会收紧信息类，处理方式相同。

\begin{proposition}[相同稀疏观测不控制完整$\KL(P\Vert Q)$ 的上界]
\label{inf:unbounded}
存在相同 Top-1 payload 的正支持参考分布族，使完整$\KL(P\Vert Q)$ 无上界，而 Tail-aware 不变。
\end{proposition}
\begin{proof}
取 $P=(0.4,0.2,0.4)$ 和 $Q_\varepsilon=(0.7,\varepsilon,0.3-\varepsilon)$，$0<\varepsilon<0.15$。全部分布的 Top-1 均为 token 1，其概率 $0.7$、尾部总质量 $0.3$ 相同。可是
\begin{equation}
\KL(P\Vert Q_\varepsilon)
=0.4\log\frac{0.4}{0.7}+0.2\log\frac{0.2}{\varepsilon}
+0.4\log\frac{0.4}{0.3-\varepsilon}\longrightarrow+\infty.
\end{equation}
第一项固定，第三项趋向有限值，第二项发散。此例需要至少两个仍未知的尾部坐标。若 payload 已查询其中一个，尾总质量会确定另一个，信息类不再保留这个自由度。
\end{proof}

\section{非线性积分与无偏性陈述的边界}
\label{inf:nonlinear}

\begin{proposition}[正的 plug-in partition 估计的 Jensen 偏差]
\label{inf:jensen}
若 $\widehat Z>0$、$\E\widehat Z=Z$ 且 $\E|\log\widehat Z|<\infty$，则
\begin{equation}
\E[\log\widehat Z]\leq\log Z.
\end{equation}
严格凹性使等号当且仅当 $\widehat Z$ 几乎处处为常数。
\end{proposition}
\begin{proof}
在 $Z>0$ 处，$\log$ 的切线满足 $\log z\leq\log Z+(z-Z)/Z$，且 $z\ne Z$ 时严格。取期望并利用 $\E(\widehat Z-Z)=0$ 即得结论；若偏离 $Z$ 的事件具有正概率，严格不等式在该事件成立，期望不等式也严格。
\end{proof}

这个命题分析 $\log\widehat Z$，并没有证明所有 $\log Z$ 估计器都不能无偏。双 token 词表中，若 anchor 全分布已知，查询任一 token 的 teacher 概率后，另一概率由归一化确定。因此一个 sampled-token 查询就能精确计算任意 $\eta$ 的 OPD partition。

更一般，设 $|\mathcal V|=D$，$m\geq D$，从已知正支持 anchor 独立抽样 $m$ 次并查询每个动作的 teacher 概率。令 $E$ 为样本覆盖整个词表的事件，$c=\Pr_q(E)>0$。$c$ 只依赖已知 anchor 与样本数。当 $E$ 成立时可恢复 teacher 全分布。定义
\begin{equation}
\widehat L=\mathbf1_E\frac{\log Z}{c},
\end{equation}
并在 $E$ 未发生时直接返回零，不调用未知 $\log Z$。逐项计算得到
\begin{equation}
\E\widehat L=\log Z,\qquad
\Var(\widehat L)=(c^{-1}-1)(\log Z)^2.
\end{equation}
该构造通常具有不可接受的成本和方差，却证明了“所有无偏 log-partition 估计都不存在”这一陈述不成立。实际方法应分别报告信息条件、估计器形式、偏差和成本。

\section{粗化、条件期望与真正的 Rao--Blackwell 化}
\label{inf:rao-blackwell}

在 $A\sim P$ 下，若知道桶内形状 KL，真实条件期望为
\begin{equation}
\E_P[\ell_A\mid C(A)=C]=\ell_C+D_C.
\end{equation}
由全方差公式，
\begin{equation}
\Var_P(\ell_A)
=\E_P[\Var_P(\ell_A\mid C(A))]
+\Var_P(\E_P[\ell_A\mid C(A)]).
\end{equation}
因此用真实条件期望替代原始样本不会增加方差。Tail-aware 仅知道 $\ell_C$，未知的 $D_C$ 没有自动出现；HETS 将头部求和并留下尾部事件修正，也不是上述条件期望。后者的方差是否下降，要由残差二阶矩直接判断。向量情形应用全协方差公式并取迹，逻辑完全相同。
